\pagestyle{fancy}
\fancyhead[l]{\autorUO}
\fancyfoot[l]{\asignaturaAbbr}
\fancyfoot[r]{\fecha}

\section{CUDA y drivers NVIDIA} \label{sec:intro_cuda}
CUDA\footnote{\url{https://developer.nvidia.com/cuda-zone}} (Compute Unified Device Architecture) es la plataforma de computación
paralela de NVIDIA\footnote{\url{https://www.nvidia.com/}} que permite ejecutar código y algoritmos en la GPU, aprovechando su
capacidad de procesamiento masivamente paralelo. \pytorch\ utiliza CUDA para acelerar las operaciones de tensores y cálculos matemáticos,
por lo que es fundamental contar con una instalación adecuada de CUDA y los drivers de NVIDIA en el sistema de desarrollo.
\subsection{Instalación de drivers NVIDIA} \label{sec:intro_nvidia_drivers}
Para instalar los drivers de NVIDIA en un sistema Linux, se pueden seguir los pasos detallados en
la documentación oficial de NVIDIA\footnote{\url{https://docs.nvidia.com/datacenter/tesla/tesla-installation-notes/index.html}}.
Es importante asegurarse de que la versión del driver sea compatible con la versión de CUDA que se va a instalar posteriormente.
\subsection{Instalación de CUDA} \label{sec:intro_cuda_install}
La instalación de CUDA puede realizarse descargando el instalador desde la página oficial de NVIDIA\
\footnote{\url{https://developer.nvidia.com/cuda-downloads}}.
Es recomendable seguir las instrucciones específicas para la distribución de Linux utilizada en el sistema de desarrollo.

\vspace{0.3cm}
\noindent{Después de la instalación, es importante verificar que CUDA y los drivers de NVIDIA funcionan correctamente con
	el comando especificado en el \autoref{lst:intro_verify_cuda}.}
\begin{lstlisting}[language=bash, caption={Verificación de la instalación de CUDA y los drivers de NVIDIA}, label={lst:intro_verify_cuda}]
nvidia-smi
\end{lstlisting}